\documentclass[11pt,a4paper]{article}

% Biên dịch bằng pdfLaTeX (chạy 2 lần để có mục lục), từ gốc repo:
%   cd docs/v3/bao_cao && pdflatex bao_cao_v3.tex && pdflatex bao_cao_v3.tex
% Các hằng cấu hình (giới hạn lượt, số ký tự...) lấy nguyên từ mã nguồn web-v2.
\usepackage[utf8]{inputenc}
\usepackage[T5]{fontenc}
\usepackage[vietnamese]{babel}
\usepackage[a4paper,left=25mm,right=22mm,top=22mm,bottom=22mm,headheight=16pt]{geometry}
\usepackage{lmodern}
\usepackage{booktabs}
\usepackage{tabularx}
\usepackage{array}
\usepackage{enumitem}
\usepackage{amsmath}
\usepackage[table]{xcolor}
\usepackage{float}
\usepackage{caption}
\usepackage{fancyhdr}
\usepackage{tikz}
\usetikzlibrary{positioning,arrows.meta,calc}
\PassOptionsToPackage{obeyspaces,spaces}{url}
\usepackage{hyperref}
\usepackage{microtype}

\definecolor{navy}{HTML}{17365D}
\definecolor{nenxanh}{HTML}{EAF2F8}
\definecolor{nenxam}{HTML}{F3F4F6}
\definecolor{nenvang}{HTML}{FFF4DC}
\definecolor{chuxam}{HTML}{666666}

\hypersetup{
  colorlinks=true, linkcolor=navy, urlcolor=navy, citecolor=navy,
  pdftitle={SafeViChi v3 - Báo cáo cải tiến - Vòng Chung kết - Chặng 2 (AI26C0007 - OneUET)},
  pdfauthor={Đội thi OneUET (AI26C0007)}
}
\setlength{\parindent}{0pt}
\setlength{\parskip}{0.5em}
\setlist{itemsep=0.2em,topsep=0.2em}
\captionsetup{font=small,labelfont=bf}
\renewcommand{\arraystretch}{1.2}
\newcolumntype{Y}{>{\raggedright\arraybackslash}X}

\newcommand{\nhan}[1]{\texttt{\detokenize{#1}}}
\newcommand{\gc}[1]{\texttt{#1}}
\newcommand{\tep}[1]{\nolinkurl{#1}}
\expandafter\def\expandafter\UrlBreaks\expandafter{\UrlBreaks\do\_\do\-\do/}
\newcommand{\ghinho}[1]{\par\noindent\colorbox{nenxanh}{\parbox{\dimexpr\linewidth-2\fboxsep}{\small #1}}\par}

\pagestyle{fancy}
\fancyhf{}
\fancyhead[L]{\small\textcolor{chuxam}{SafeViChi v3}}
\fancyhead[R]{\small\textcolor{chuxam}{Đội OneUET $\cdot$ Mã AI26C0007}}
\fancyfoot[L]{\footnotesize\textcolor{chuxam}{Cuộc thi Sáng tạo trẻ Quốc gia trong lĩnh vực Trí tuệ nhân tạo năm 2026 $\cdot$ Vòng Chung kết -- Chặng 2}}
\fancyfoot[R]{\small\textcolor{navy}{\textbf{\thepage}}}
\renewcommand{\headrulewidth}{0.4pt}
\renewcommand{\footrulewidth}{0.4pt}
\renewcommand{\headrule}{{\color{navy!35}\hrule height \headrulewidth width \headwidth \vskip-\headrulewidth}}
\renewcommand{\footrule}{\vskip-\footruleskip\vskip-\footrulewidth{\color{navy!25}\hrule height \footrulewidth width \headwidth}\vskip\footruleskip}

\begin{document}

% ---------- Trang bìa ----------
\hypersetup{pageanchor=false}
\begin{titlepage}
\centering
\vspace*{18mm}
{\large\bfseries\color{navy} CUỘC THI SÁNG TẠO TRẺ QUỐC GIA\\[2pt] TRONG LĨNH VỰC TRÍ TUỆ NHÂN TẠO NĂM 2026\par}
\vspace{4mm}
{\normalsize\bfseries\color{chuxam} VÒNG CHUNG KẾT -- CHẶNG 2\par}

\vspace{34mm}
{\Large Báo cáo cải tiến\par}
\vspace{6mm}
{\fontsize{34}{40}\selectfont\bfseries\color{navy} SafeViChi v3\par}
\vspace{6mm}
{\Large Cơ chế phản hồi, kháng nghị và hiệu chỉnh nhãn\\[2pt] do con người quyết định, tối thiểu hóa dữ liệu cá nhân\par}

\vspace{26mm}
{\renewcommand{\arraystretch}{1.35}
\begin{tabular}{@{}l@{\hspace{14mm}}l@{}}
\textbf{\color{chuxam}Đội thi} & \textbf{OneUET} \\
\textbf{\color{chuxam}Mã đội thi} & \textbf{\texttt{AI26C0007}} \\
\end{tabular}}
\vfill
\end{titlepage}
\hypersetup{pageanchor=true}

\tableofcontents
\newpage

% =====================================================================
\section{Tổng quan}

\ghinho{\textbf{Phiên bản v3 giữ nguyên mô hình của v2.} Bộ nhãn, mô hình \nhan{SafeVi-DualORC-CRF} và cơ chế
phân tích tin nhắn trên thiết bị không thay đổi so với bản đã nộp ở Chặng 1. Báo cáo này chỉ trình bày
các thành phần \textbf{bổ sung}: cơ chế để người dùng phản hồi khi kết quả của mô hình chưa đúng, và quy trình
để người kiểm duyệt xem xét, hiệu chỉnh nhãn. Mọi quyết định cuối cùng đều do con người đưa ra.}

Bảng~\ref{tab:debai} tóm tắt cách v3 đáp ứng từng yêu cầu của đề bài.

\begin{table}[H]
\centering
\small
\begin{tabularx}{\linewidth}{@{}p{4.6cm}Y@{}}
\toprule
\textbf{Yêu cầu của đề bài} & \textbf{Giải pháp của v3} \\
\midrule
Gán nhãn độc hại tự động có thể ảnh hưởng đến quyền tự do ngôn luận và phát sinh khiếu nại &
Mô hình chỉ \textbf{cảnh báo}, không chặn và không xóa tin nhắn; nút \textbf{Vẫn gửi} luôn hiển thị.
Người dùng có thể phản hồi ngay trên kết quả phân tích (mục~\ref{sec:phanhoi}). \\
\addlinespace
Cần cơ chế phản hồi, kháng nghị, hiệu chỉnh nhãn và giới hạn quyền quyết định cuối cùng của AI &
Phản hồi có đính kèm đoạn hội thoại được đưa vào \textbf{hàng đợi kiểm duyệt}; \textbf{người kiểm duyệt}
xem xét từng mẫu và xác định nhãn đúng cùng bằng chứng (mục~\ref{sec:duyet}).
Mô hình không tự thay đổi nhãn. Mọi thao tác đều được ghi nhật ký. \\
\addlinespace
Giảm lưu trữ thông tin nhạy cảm &
Phản hồi chỉ gồm các lựa chọn định sẵn, \textbf{không chứa nội dung tin nhắn}. Đoạn hội thoại chỉ được đính
kèm khi người dùng chủ động chọn và xác nhận đồng ý, sau khi \textbf{bộ lọc thông tin nhận dạng cá nhân}
đã che thông tin trên thiết bị (mục~\ref{sec:dinhkem}, \ref{sec:dulieu}). \\
\bottomrule
\end{tabularx}
\caption{Đối chiếu yêu cầu của đề bài với giải pháp của v3.}
\label{tab:debai}
\end{table}

\newpage
% =====================================================================
\section{Quy trình xử lý tổng thể}
\label{sec:pipeline}

Hình~\ref{fig:pipeline} mô tả toàn bộ quy trình của v3. Phần phân tích tin nhắn vẫn chạy hoàn toàn trên
thiết bị như ở v2. Dữ liệu chỉ được gửi lên máy chủ khi người dùng \textbf{chủ động gửi phản hồi}.

\begin{figure}[H]
\centering
\begin{tikzpicture}[
  font=\small,
  hop/.style={draw=navy, rounded corners=3pt, fill=nenxanh, align=center, inner sep=4pt, text width=6.1cm},
  rong/.style={hop, text width=13.4cm},
  may/.style={hop, fill=white},
  nguoi/.style={hop, fill=nenvang},
  mui/.style={-{Stealth[length=2.2mm]}, navy, thick},
  >=Stealth]
% Trên thiết bị
\node[rong] (a) at (0,0) {Tin nhắn $\to$ mô hình v2 (ONNX, chạy trên thiết bị) $\to$ nhãn và cụm bằng chứng
  \\\textcolor{chuxam}{\footnotesize không phát sinh yêu cầu mạng khi phân tích}};
\node[rong] (b) at (0,-1.35) {\textbf{Cảnh báo}: mô hình chỉ gợi ý, không chặn, không xóa; nút \textbf{Vẫn gửi} luôn hiển thị};
\node[rong, fill=nenvang] (c) at (0,-2.6) {Người dùng chọn \textbf{Phản hồi kết quả} hoặc \textbf{Báo AI hiểu sai}};
\node[rong] (fb) at (0,-3.85) {\textbf{Phản hồi}: chọn loại báo cáo và lý do từ các lựa chọn định sẵn
  \\\textcolor{chuxam}{\footnotesize không chứa nội dung tin nhắn}};
\node[hop, text width=4.2cm, fill=white] (dec) at (0,-5.15) {Đính kèm đoạn hội thoại?\\
  \textcolor{chuxam}{\footnotesize hỏi ở bước cuối\\ của mỗi lần phản hồi}};
\node[may] (m1) at (-3.65,-6.65) {\textbf{Không, chỉ gửi phản hồi}\\ \textcolor{chuxam}{\footnotesize chỉ gồm các lựa chọn}};
\node[nguoi] (m2) at (3.65,-6.65) {\textbf{Đồng ý gửi kèm đoạn hội thoại}\\ bộ lọc che thông tin cá nhân\\ trên thiết bị};
% Máy chủ
\node[rong, fill=nenxam] (api) at (0,-8.3) {\textbf{API tiếp nhận phản hồi} (cùng nguồn gốc với ứng dụng web):\\
  kiểm tra định dạng, chống gửi trùng, giới hạn số lượt mỗi ngày $\to$ lưu trữ};
\node[may] (t1) at (-3.65,-9.8) {Thống kê mọi phản hồi\\ (chỉ đếm số lượt)};
\node[nguoi, below=5mm of t1] (t2) {Người kiểm duyệt xem các cặp nhãn\\ và lý do được phản hồi nhiều};
\node[may] (r1) at (3.65,-9.8) {Hàng đợi kiểm duyệt (bản đã che)};
\node[nguoi, below=5mm of r1] (r2) {\textbf{Người kiểm duyệt}\\ che thêm thông tin còn sót,\\ xác định nhãn và bằng chứng};
\node[may, below=5mm of r2] (r3) {Nhãn do con người gán\\ \nhan{nguon=phan_hoi}, \nhan{cach_gan=nguoi}};
\node[may, below=5mm of r3] (r4) {Chỉ đưa vào \textbf{tập huấn luyện},\\ không đưa vào tập kiểm thử};
\path (r4.south) ++(0,-0.6) coordinate (nkpos);
\node[rong, draw=chuxam, dashed, fill=white, anchor=north] (nk) at (nkpos -| a) {\textbf{Nhật ký}: mở mẫu, che thêm, lưu kết quả kiểm duyệt\\
  (người thực hiện, thời điểm, thao tác)};
\draw[mui] (a) -- (b);
\draw[mui] (b) -- (c);
\draw[mui] (c) -- (fb);
\draw[mui] (fb) -- (dec);
\draw[mui] (dec.west) -| node[pos=0.25, above, font=\footnotesize]{Không} (m1.north);
\draw[mui] (dec.east) -| node[pos=0.25, above, font=\footnotesize]{Có} (m2.north);
\draw[mui] (m1.south) -- (m1.south |- api.north);
\draw[mui] (m2.south) -- (m2.south |- api.north);
\draw[mui] (api.south -| t1.north) -- (t1.north);
\draw[mui] (api.south -| r1.north) -- (r1.north);
\draw[mui] (t1) -- (t2);
\draw[mui] (r1) -- (r2); \draw[mui] (r2) -- (r3); \draw[mui] (r3) -- (r4);
\node[font=\footnotesize\bfseries, text=chuxam, anchor=south west] at (-6.9,0.45) {TRÊN THIẾT BỊ NGƯỜI DÙNG};
\node[font=\footnotesize\bfseries, text=chuxam, anchor=south west] at (-6.9,-7.8) {MÁY CHỦ};
\end{tikzpicture}
\caption{Quy trình xử lý của v3. Ô màu vàng là bước do \textbf{con người quyết định}; mô hình không có bước
nào tự thay đổi nhãn.}
\label{fig:pipeline}
\end{figure}

Quy trình gồm năm bước:
\begin{enumerate}[leftmargin=*]
\item \textbf{Phân tích và cảnh báo} (giữ nguyên v2): mô hình chạy trên thiết bị, trả về nhãn và cụm bằng chứng.
  Người dùng luôn có thể chọn \textbf{Vẫn gửi}.
\item \textbf{Phản hồi}: khi cho rằng kết quả chưa đúng, người dùng gửi phản hồi bằng các lựa chọn định sẵn.
  Nếu cần người kiểm duyệt xem xét ngữ cảnh, người dùng có thể đính kèm đoạn hội thoại đã che.
\item \textbf{Tiếp nhận}: máy chủ kiểm tra từng trường dữ liệu và từ chối gói tin không hợp lệ trước khi lưu.
\item \textbf{Xử lý theo hai nhánh}: mọi phản hồi được tổng hợp vào \textbf{thống kê}; phản hồi có đính
  kèm được đưa thêm vào \textbf{hàng đợi kiểm duyệt}.
\item \textbf{Không tự học}: nhãn do con người xác định chỉ được dùng cho lần huấn luyện tiếp theo;
  mô hình đang vận hành không tự thay đổi theo phản hồi.
\end{enumerate}

% =====================================================================
\section{Cơ chế phản hồi kết quả}
\label{sec:phanhoi}

Phản hồi là một thao tác duy nhất. Mọi phản hồi gồm các lựa chọn định sẵn và không chứa nội dung tin nhắn
(mục~\ref{sec:noidung}). Khi cần người kiểm duyệt xem xét ngữ cảnh, người dùng có thể \textbf{đính kèm đoạn
hội thoại đã che} (mục~\ref{sec:dinhkem}); khi đó phản hồi được đưa vào hàng đợi kiểm duyệt. Phản hồi không
đính kèm chỉ được dùng để thống kê.

Biểu mẫu gồm hai bước. Bước thứ nhất hỏi các lựa chọn của phản hồi; bước cuối, \textbf{Chia sẻ đoạn hội thoại?},
hỏi người dùng có đồng ý gửi kèm đoạn hội thoại hay không. Bước cuối được hỏi lại ở \textbf{mỗi lần phản hồi},
hệ thống không ghi nhớ lựa chọn của lần trước.

\subsection{Nội dung phản hồi}
\label{sec:noidung}

\textbf{Điểm truy cập.} Nút \textbf{Phản hồi kết quả} nằm trên dải kết quả phân tích; nút
\textbf{Báo AI hiểu sai} nằm trong cảnh báo trước khi gửi tin. Người dùng chọn một trong năm loại báo cáo
(bảng~\ref{tab:loaibaocao}); mỗi loại có câu hỏi bổ sung riêng, đều ở dạng lựa chọn định sẵn.

\begin{table}[H]
\centering
\small
\begin{tabularx}{\linewidth}{@{}p{4.6cm}p{3.1cm}Y@{}}
\toprule
\textbf{Loại báo cáo} & \textbf{Điều kiện hiển thị} & \textbf{Câu hỏi bổ sung (ví dụ)} \\
\midrule
Tin nhắn này an toàn & Mô hình gắn nhãn nguy hại & Lý do: nói đùa, cách nói thân mật, từ lóng, phương ngữ, trích dẫn, ngữ cảnh thông thường \\
Có vấn đề, nhưng AI chọn sai loại & Mô hình gắn nhãn nguy hại & Nhãn phù hợp hơn và dấu hiệu nhận biết \\
AI bỏ sót nội dung nguy hại & Mô hình gắn nhãn An toàn & Loại nội dung; dấu hiệu trực tiếp, các lượt trước hoặc cách viết biến dạng \\
Nhãn đúng, phần tô sáng có vấn đề & Mô hình gắn nhãn nguy hại & Tô sáng nhầm, tô sáng thiếu hoặc chưa tô sáng \\
Vấn đề khác hoặc tôi chưa chắc & Luôn hiển thị & Điểm còn phân vân; cảnh báo khó hiểu hoặc xuất hiện quá thường xuyên \\
\bottomrule
\end{tabularx}
\caption{Năm loại báo cáo trong biểu mẫu phản hồi.}
\label{tab:loaibaocao}
\end{table}

\textbf{Thành phần gói tin.} Gói tin chỉ gồm các trường dưới đây; máy chủ từ chối mọi trường khác.
\begin{itemize}[leftmargin=*]
\item Mã loại báo cáo, mã lý do và mã chi tiết, đều là mã cố định (ví dụ \nhan{tin_nhan_an_toan},
  \nhan{dua_noi_than_mat}).
\item Nhãn của mô hình, nhãn do người dùng đề xuất (được để trống khi người dùng chưa chắc chắn) và trạng thái
  có tô sáng bằng chứng hay không.
\item \textbf{Khoảng điểm} từ 0 đến 9 (xác suất làm tròn xuống theo bước 0,1), phiên bản mô hình và phiên bản ngưỡng.
\item Mã lần gửi để chống gửi trùng và mã thiết bị; cả hai là mã định danh ngẫu nhiên (UUID), mã thiết bị được
  lưu trong bộ nhớ cục bộ của trình duyệt.
\item \textbf{Mô tả ngắn}, chỉ có khi người dùng chọn ``Khác'': tối đa 200 ký tự, đã được bộ lọc thông tin nhận
  dạng cá nhân xử lý trên thiết bị; máy chủ chạy lại bộ lọc để kiểm tra.
\end{itemize}

\textbf{Dữ liệu không được gửi:} nội dung tin nhắn, cụm bằng chứng, mã tài khoản đăng nhập, số điện thoại và
tên liên hệ. Máy chủ \textbf{chỉ ghi ngày} tiếp nhận, không ghi giờ phút. Mã thiết bị \textbf{không được lưu cùng
phản hồi}; mã này chỉ dùng để giới hạn tối đa 20 phản hồi không đính kèm và 5 phản hồi có đính kèm mỗi ngày,
dữ liệu đếm của các ngày trước được xóa.

\textbf{Thống kê phản hồi.} Mọi phản hồi, có hoặc không đính kèm, đều được tổng hợp trong mục
\textbf{Thống kê phản hồi}: tổng số phản hồi, số phản hồi có đính kèm (đã duyệt và chờ duyệt), số lượt theo loại báo
cáo, các cặp nhãn gặp nhiều nhất (nhãn của mô hình và nhãn người dùng đề xuất) và các lý do phổ biến. Đây là dữ liệu do người dùng tự báo cáo, vì vậy chỉ
dùng để \textbf{xác định các trường hợp cần xem xét}, không dùng để ước lượng tỷ lệ lỗi của mô hình và không
được đưa vào huấn luyện.

\subsection{Đính kèm đoạn hội thoại}
\label{sec:dinhkem}

Tùy chọn đính kèm dành cho trường hợp người dùng muốn \textbf{kháng nghị} và cần người kiểm duyệt xem xét ngữ cảnh.
Phản hồi giữ nguyên các lựa chọn ở mục~\ref{sec:noidung} và bổ sung đoạn hội thoại đã che. Quy trình gồm bốn bước;
việc che thông tin được thực hiện trên thiết bị trước khi dữ liệu được gửi đi:

\begin{enumerate}[leftmargin=*]
\item \textbf{Lọc tự động.} Khi người dùng chuyển sang bước cuối, \textbf{bộ lọc thông tin nhận dạng cá nhân}
  thay thông tin cá nhân trong đoạn hội thoại (tối đa 8 lượt) bằng nhãn giữ chỗ; quy tắc chi tiết được trình bày
  ở mục~\ref{sec:xuly}.
\item \textbf{Thông báo và cam kết.} Bước \textbf{Chia sẻ đoạn hội thoại?} giải thích mục đích của việc gửi kèm,
  nêu rõ việc chia sẻ là \textbf{hoàn toàn tự nguyện}, và trình bày cam kết bảo vệ quyền riêng tư: thông tin nhận
  dạng được che trên thiết bị; không đính kèm mã tài khoản, số điện thoại hay địa chỉ IP; chỉ đội ngũ kiểm duyệt
  được truy cập nội dung; dữ liệu không được mua bán, chia sẻ với bên thứ ba hoặc dùng cho mục đích thương mại.
\item \textbf{Người dùng quyết định.} Người dùng chọn một trong hai nút: \textbf{Đồng ý gửi kèm đoạn hội thoại}
  hoặc \textbf{Không, chỉ gửi phản hồi}; nút \textbf{Quay lại} cho phép sửa các lựa chọn trước đó. Không có lựa
  chọn nào được chọn sẵn.
\item \textbf{Gửi phản hồi.} Khi người dùng đồng ý, gói tin gồm các trường của phản hồi, bản đoạn hội thoại đã che
  và phiên bản nội dung cam kết mà người dùng đã đồng ý. Máy chủ chạy lại bộ lọc để bảo đảm không còn thông tin
  nhận dạng được, kiểm tra đoạn hội thoại không quá 8 lượt và giới hạn tối đa 5 phản hồi có đính kèm mỗi ngày.
\end{enumerate}

\newpage
% =====================================================================
\section{Quy trình xử lý dữ liệu khi gửi}
\label{sec:xuly}

Dữ liệu phản hồi đi qua hai giai đoạn: \textbf{xử lý trên thiết bị} trước khi rời khỏi trình duyệt và
\textbf{kiểm tra trên máy chủ} trước khi được lưu (hình~\ref{fig:xuly}).

\begin{figure}[H]
\centering
\begin{tikzpicture}[font=\small, node distance=4mm,
  bk/.style={draw=navy, rounded corners=3pt, fill=nenxanh, align=center, inner sep=4pt, text width=2.55cm, minimum height=1.45cm},
  sv/.style={bk, fill=nenxam},
  mui/.style={-{Stealth[length=2mm]}, navy, thick}]
\node[bk] (s1) {Lấy tối đa\\ 8 lượt gần nhất};
\node[bk, right=of s1] (s2) {Chuẩn hóa\\ ký tự};
\node[bk, right=of s2] (s3) {Bộ lọc thông tin\\ nhận dạng cá nhân};
\node[bk, right=of s3] (s4) {Đóng gói\\ phản hồi};
\node[sv, below=9mm of s4] (v1) {Kiểm tra\\ định dạng};
\node[sv, left=of v1] (v2) {Chạy lại\\ bộ lọc};
\node[sv, left=of v2] (v3) {Giới hạn\\ số lượt};
\node[sv, left=of v3] (v4) {Lưu trữ\\ tối thiểu};
\draw[mui] (s1) -- (s2); \draw[mui] (s2) -- (s3); \draw[mui] (s3) -- (s4);
\draw[mui] (s4) -- (v1); \draw[mui] (v1) -- (v2); \draw[mui] (v2) -- (v3); \draw[mui] (v3) -- (v4);
\node[font=\footnotesize\bfseries, text=chuxam, anchor=south west] at (s1.north west) {TRÊN THIẾT BỊ};
\node[font=\footnotesize\bfseries, text=chuxam, anchor=north west] at (v4.south west) {TRÊN MÁY CHỦ};
\end{tikzpicture}
\caption{Các bước xử lý dữ liệu phản hồi.}
\label{fig:xuly}
\end{figure}

\subsection{Xử lý trên thiết bị}

\begin{enumerate}[leftmargin=*]
\item \textbf{Lấy đoạn hội thoại.} Ứng dụng lấy đúng các lượt mô hình vừa phân tích (tối đa 8 lượt gần nhất).
  Mỗi lượt là một dòng; ký tự xuống dòng bên trong một lượt được đổi thành dấu cách.
\item \textbf{Chuẩn hóa ký tự.} Văn bản được chuẩn hóa Unicode (NFKC) và loại bỏ các ký tự vô hình (ký tự độ rộng
  bằng không, ký tự điều khiển hướng chữ) để thông tin cá nhân không lọt qua bộ lọc nhờ các ký tự này.
\item \textbf{Lọc thông tin nhận dạng cá nhân.} Bộ lọc nhận diện và thay thế thông tin theo bảng~\ref{tab:boloc}.
  Việc nhận diện từ khóa được thực hiện trên bản đã bỏ dấu, nên các cách viết như ``sdt'', ``so tk'', ``tai khoan''
  vẫn được nhận ra; chuỗi số viết tách bằng dấu cách, dấu chấm hoặc dấu gạch vẫn được xem là một số.
\item \textbf{Đóng gói.} Gói tin gồm các lựa chọn của phản hồi, mô tả ngắn (nếu có, đã qua cùng bộ lọc) và, khi người
  dùng đồng ý, bản đoạn hội thoại đã che. Yêu cầu được gửi \textbf{không kèm cookie hay thông tin đăng nhập}.
\end{enumerate}

\begin{table}[!htbp]
\centering
\footnotesize
\renewcommand{\arraystretch}{1.1}
\begin{tabularx}{\linewidth}{@{}p{2.6cm}Yp{2.2cm}@{}}
\toprule
\textbf{Loại thông tin} & \textbf{Cách nhận diện} & \textbf{Thay bằng} \\
\midrule
Đường link & Chuỗi bắt đầu bằng \texttt{http://}, \texttt{https://}, \texttt{www.} hoặc tên miền có đường dẫn, tham số;
  link chứa tên đăng nhập, mật khẩu cũng được che toàn bộ & \gc{[LINK]} \\
Email & Chuỗi có dạng \texttt{tên@tên-miền} & \gc{[EMAIL]} \\
Mã OTP & Chuỗi 4--8 chữ số đứng sau từ khóa ``OTP'', ``mã'', ``code'', ``PIN'', ``mã xác minh'' & \gc{[MÃ]} \\
Giấy tờ tùy thân & Chuỗi 9 hoặc 12 chữ số sau ``CCCD'', ``CMND'', ``căn cước'', ``chứng minh''; hoặc chuỗi 9, 12 chữ
  số không có từ khóa & \gc{[GIẤY\_TỜ]} \\
Số tài khoản, số thẻ & Chuỗi 6--19 chữ số sau ``STK'', ``tài khoản'', ``ngân hàng'', ``ATM'' hoặc tên ngân hàng, ví
  điện tử (Vietcombank, Techcombank, MoMo, ZaloPay\ldots); hoặc chuỗi 10--19 chữ số không có từ khóa & \gc{[STK]} \\
Số điện thoại & Chuỗi 9--12 chữ số sau ``SĐT'', ``điện thoại'', ``Zalo'', ``hotline'', ``liên hệ'', ``gọi'';
  hoặc số bắt đầu bằng 0 có 10--11 chữ số, số dạng \texttt{+84} & \gc{[SĐT]} \\
Địa chỉ & Số nhà đi kèm các thành phần đường, phố, ngõ, ngách, hẻm, phường, xã, quận, huyện, tỉnh, thành phố,
  thôn, ấp\ldots & \gc{[ĐỊA\_CHỈ]} \\
Tên người & Tên liên hệ và tên người tham gia cuộc trò chuyện; tên được nhắc bằng \texttt{@}; cụm viết hoa đứng
  sau từ xưng hô hoặc giới thiệu (anh, chị, em, cô, chú, ``tên là'', ``mình là''\ldots), trừ thán từ, từ ngữ
  thông dụng và tên ngân hàng, dịch vụ; một tên đã nhận diện xuất hiện lại ở chỗ khác cũng được che &
  \gc{[TÊN\_1]}, \gc{[TÊN\_2]}\ldots \\
\midrule
Số tiền & Số đi kèm đơn vị (k, tr, triệu, nghìn, đ, VND, đồng, USD, tỷ, \%\ldots) & \textbf{Giữ nguyên} \\
\bottomrule
\end{tabularx}
\caption{Quy tắc của bộ lọc thông tin nhận dạng cá nhân. Chuỗi số được xét từ khóa trong phạm vi 30 ký tự đứng
trước, trên cùng một dòng.}
\label{tab:boloc}
\end{table}

Số tiền được giữ nguyên vì đây là căn cứ quan trọng để người kiểm duyệt nhận diện lừa đảo; vùng số tiền được đánh
dấu trước nên không bị nhầm thành số tài khoản hay số điện thoại. Các tên giống nhau nhận cùng một số thứ tự
(ví dụ ``Nguyễn Thị Lan'' và lần nhắc sau đó ``Lan'' cùng là \gc{[TÊN\_1]}) để người kiểm duyệt vẫn theo dõi được ai nói với ai.
Bảng đánh số này chỉ tồn tại trong lúc bộ lọc chạy và không được lưu. Tên liên hệ chỉ dùng để nhận diện trên
thiết bị, không được gửi đi. Đây là phép \textbf{che, ẩn danh hóa một chiều}, không phải mã hóa: không có bảng
đối chiếu với thông tin thật và không thể khôi phục bản gốc.

\subsection{Kiểm tra trên máy chủ}

\begin{enumerate}[leftmargin=*]
\item \textbf{Kiểm tra định dạng.} Máy chủ chỉ chấp nhận các trường trong danh sách cho phép; kiểm tra kiểu dữ
  liệu, mã lựa chọn và tổ hợp lựa chọn hợp lệ. Gói tin không hợp lệ bị từ chối và không được lưu.
\item \textbf{Chạy lại bộ lọc.} Máy chủ chạy lại cùng bộ lọc trên mô tả ngắn và đoạn hội thoại; nếu bộ lọc còn
  phát hiện thông tin cần che, gói tin bị từ chối. Đoạn hội thoại phải có tối đa 8 lượt và không có lượt rỗng.
\item \textbf{Giới hạn số lượt.} Mỗi thiết bị gửi tối đa 20 phản hồi không đính kèm và 5 phản hồi có đính kèm mỗi
  ngày. Gửi lại cùng một phản hồi (cùng mã lần gửi) không tạo bản ghi mới.
\item \textbf{Lưu trữ tối thiểu.} Máy chủ lưu nội dung phản hồi kèm \textbf{ngày} tiếp nhận (không lưu giờ phút).
  Mã thiết bị không được lưu cùng phản hồi; dữ liệu đếm số lượt của các ngày trước được xóa.
\item \textbf{Truy cập sau khi lưu.} Đoạn hội thoại chỉ được người kiểm duyệt mở theo từng mẫu, mỗi lần mở đều ghi
  nhật ký; nếu phát hiện thông tin còn sót, người kiểm duyệt che thêm trước khi duyệt (mục~\ref{sec:duyet}).
\end{enumerate}

% =====================================================================
\section{Kiểm duyệt và hiệu chỉnh nhãn}
\label{sec:duyet}

Khu vực kiểm duyệt là tab \textbf{Kiểm duyệt \& quản trị} nằm cạnh tab \textbf{Nhắn tin \& cảnh báo} trong ứng
dụng; bản trình diễn có một vai trò là \textbf{Người kiểm duyệt}.

\begin{itemize}[leftmargin=*]
\item \textbf{Hàng đợi kiểm duyệt} chỉ gồm phản hồi có đính kèm, chia thành hai mục \textbf{Chờ duyệt} và
  \textbf{Đã duyệt}. Danh sách chỉ hiển thị thông tin tóm tắt (ngày, nhãn của mô hình, nhãn đề xuất, loại báo
  cáo, lý do). Đoạn hội thoại chỉ được xem khi mở từng mẫu, và mỗi lần mở đều được \textbf{ghi nhật ký}.
  Hệ thống không cung cấp API xuất dữ liệu hàng loạt.
\item Mỗi mẫu hiển thị ba phần: \textbf{Kết quả AI} (nhãn, khoảng độ tin cậy, có tô sáng bằng chứng hay không;
  bản thân cụm bằng chứng của mô hình không được gửi lên máy chủ), \textbf{Người dùng báo} (nhãn đề xuất, loại báo
  cáo, lý do, mô tả) và \textbf{Quyết định của bạn}.
\item Nếu bản đã che còn sót thông tin cá nhân, người kiểm duyệt chọn phần văn bản và bấm \textbf{Che phần đã chọn}
  để thay bằng \gc{[ĐÃ\_CHE]}. Thao tác này chỉ thực hiện được khi mẫu chưa có lượt duyệt nào và được ghi nhật ký.
\item Người kiểm duyệt gán nhãn theo hướng dẫn gán nhãn đã dùng ở v2: chọn \textbf{nhãn sau khi xem xét}, đánh dấu
  nếu nội dung nhắm vào một nhóm hoặc vùng miền, chọn \textbf{cụm bằng chứng} trực tiếp trên bản đã che (nhờ đó
  bằng chứng luôn là trích dẫn nguyên văn) và có thể thêm ghi chú ngắn (tối đa 200 ký tự).
\item Mỗi người chỉ kiểm duyệt một mẫu một lần; \textbf{một lượt duyệt} chuyển mẫu từ \nhan{cho_duyet} sang
  \nhan{da_duyet}, và mẫu đã duyệt không thể duyệt lại.
\end{itemize}

Khi đạt trạng thái \nhan{da_duyet}, mẫu được gắn \nhan{nguon="phan_hoi"}, \nhan{cach_gan="nguoi"} và \textbf{chỉ
được đưa vào tập huấn luyện} của lần huấn luyện tiếp theo. Việc xuất mẫu đã duyệt được thực hiện thủ công, ngoài
ứng dụng, và tự động loại các câu trùng nguyên văn với tập kiểm thử. Mẫu phản hồi không bao giờ được đưa vào tập kiểm thử đã cố định, và hệ
thống không có cơ chế tự học: mô hình v2 hiện tại giữ nguyên.

% =====================================================================
\section{Phân định quyền quyết định}
\label{sec:quyen}

\begin{table}[H]
\centering
\small
\begin{tabularx}{\linewidth}{@{}p{3.3cm}YY@{}}
\toprule
\textbf{Chủ thể} & \textbf{Được phép} & \textbf{Không được phép} \\
\midrule
Mô hình (AI) & Gợi ý nhãn, tô sáng bằng chứng, hiển thị cảnh báo & Chặn gửi, xóa tin nhắn, tự thay đổi nhãn, tự học từ phản hồi \\
Người dùng & Chọn \textbf{Vẫn gửi}; gửi phản hồi; quyết định ở mỗi lần phản hồi có đính kèm đoạn hội thoại hay không & Quyết định nhãn đưa vào dữ liệu huấn luyện \\
Người kiểm duyệt & Xác định nhãn và bằng chứng cho phản hồi có đính kèm; che thêm thông tin còn sót; xem thống kê phản hồi & Duyệt lại mẫu đã chốt; che thêm sau khi mẫu đã có lượt duyệt \\
\bottomrule
\end{tabularx}
\caption{Quyền quyết định của từng chủ thể.}
\end{table}

\textbf{Điều chỉnh mức cảnh báo theo người dùng.} Ứng dụng ghi nhận trên thiết bị số lần người dùng báo
``Tin nhắn này an toàn'' cho từng nhãn. Khi đạt 3 lần trong 7 ngày, cách hiển thị cảnh báo được điều chỉnh:
\nhan{XUC_PHAM} không còn làm mờ tin nhắn, \nhan{QUAY_ROI} không còn gợi ý chặn. Hai nhãn \nhan{DE_DOA} và
\nhan{LUA_DAO} có \textbf{mức cảnh báo tối thiểu} và luôn giữ nguyên cảnh báo. Dữ liệu đếm chỉ lưu trên thiết bị;
mô hình không thay đổi.

% =====================================================================
\section{Tối thiểu hóa dữ liệu cá nhân}
\label{sec:dulieu}

\begin{table}[H]
\centering
\small
\begin{tabularx}{\linewidth}{@{}p{4.4cm}p{2.8cm}Y@{}}
\toprule
\textbf{Dữ liệu} & \textbf{Gửi lên máy chủ} & \textbf{Ghi chú} \\
\midrule
Tin nhắn khi phân tích & Không & Như v2, không phát sinh yêu cầu mạng \\
Phản hồi: mã lựa chọn, nhãn, khoảng điểm & Khi người dùng gửi & Chỉ lưu ngày, không lưu giờ phút \\
Mô tả ``Khác'' & Khi người dùng gửi & Đã qua bộ lọc thông tin nhận dạng cá nhân, tối đa 200 ký tự \\
Đoạn hội thoại đính kèm & Chỉ khi người dùng đồng ý & Đã che trên thiết bị; hỏi lại ở mỗi lần phản hồi \\
Mã thiết bị & Khi người dùng gửi & Chỉ dùng để giới hạn số lượt mỗi ngày; không lưu cùng phản hồi \\
Tài khoản đăng nhập, số điện thoại, tên liên hệ & Không & Không có trong gói tin \\
Dữ liệu đếm để điều chỉnh cảnh báo & Không & Chỉ lưu trong bộ nhớ cục bộ của trình duyệt \\
\bottomrule
\end{tabularx}
\caption{Dữ liệu được gửi và lưu trữ trong v3.}
\end{table}

Máy chủ của bản trình diễn hoạt động cùng nguồn gốc với ứng dụng web (Cloudflare Pages Functions và cơ sở dữ liệu
Cloudflare D1), vì vậy không cần mở rộng chính sách bảo mật nội dung của trang. Cơ sở dữ liệu của bản trình diễn
chỉ chứa dữ liệu giả lập.

\end{document}
